\documentclass[10pt,a4paper]{amsart}
\usepackage[margin=19mm]{geometry}
\usepackage{amsmath,amssymb,mathtools}
\usepackage[scr=rsfs,cal=euler]{mathalfa}
\usepackage{xfrac}
\usepackage[unicode,hidelinks,bookmarks=false]{hyperref}
\newcommand{\ie}{i.e.\ }
\newcommand{\cf}{cf.\ }
\newcommand{\resp}{respectively\ }
\newcommand{\Subsection}{\subsection}
\providecommand{\nobreakdash}{\nobreak\hbox{-}}
\setlength{\emergencystretch}{2em}
\multlinegap0pt
\usepackage{mathtools}
\usepackage{tikz}
\usepackage[shortlabels]{enumitem}
\usepackage[normalem]{ulem}
\newtheorem{lemma}{Lemma}
\title{A new antisymmetrizer-to-determinant formula and a conjecture of Colomo and Pronko}
\author{Ilse Fischer, Markus Reibnegger}
\date{Source: arXiv:2608.24548v1 (CC BY 4.0)}
\begin{document}
\maketitle

\noindent\emph{Source and licence.} A new antisymmetrizer-to-determinant formula and a conjecture of Colomo and Pronko. Version arXiv:2608.24548v1 (CC BY 4.0), distributed by arXiv under the Creative Commons Attribution 4.0 International licence, http://creativecommons.org/licenses/by/4.0/, as recorded in the arXiv licence metadata for this exact version and shown on the abstract page https://arxiv.org/abs/2608.24548v1. Full licence notice: \texttt{licenses/arxiv-2608.24548-CC-BY-4.0.txt}.\par\smallskip

\noindent\textit{Editorial scope.} Counted unit: the article's lemma generalcoeff (source0.tex lines 715--735). The article's complete original proof of it is delivered with it (source0.tex lines 737--808, one \texttt{proof} environment of 72 lines). No mathematics is written, repaired or re-derived: the statement and the proof are the article's own lines, in the article's own notation, and no retained line is substituted or re-flowed.  Retained uncounted as the source-local context the proof needs: lines 564--579.  Nothing was omitted from any retained span, so the package carries no omission record.  No retained line contains a citation, so the wrapper carries no bibliography and no bibliography entry.  No retained line contains a figure environment, an image-inclusion command or drawing code, so no image is delivered and the package declares no asset dependency.  Editorial formatting only, in addition: an AMS \texttt{amsart} preamble that restates the article's own declarations and macros used by the retained lines, namely the environments \texttt{lemma} on separate counters, together with the standard packages those lines need.  The retained environments print in this document as Lemma 1 in order of appearance, whereas the article prints the same items with the numbering of its own full document.  The complete native source of the article as distributed in the arXiv e-print for this exact version is bundled unchanged as \texttt{sources/208489/source0.tex}.
\begin{lemma}
\label{coeffbase}
For any positive even integer $n$ and any $i \in [n]$, we have 
$$
\sum_{j=1}^n \sum_{k=0}^n a_{j,k} e_k(q(X_1),\ldots,q(X_n))p(X_i)^{\frac{n}{2}}(q(X_i)^{-j}-p(X_i)^{-j}) = 
\prod_{k=1}^n (p(X_i) - q(X_k)),
$$ 
where 
$$
a_{j,k} = 
\begin{cases} (-1)^k \left( \binom{\frac{n}{2}+j-k}{n-2k-j} + \binom{\frac{n}{2}+j-k-1}{n-2k-j-1}\right), & \text{if 
$k \le \frac{n}{2}$}, \\
                       (-1)^{k+1} \left[k=\frac{n}{2}+j\right], & \text{otherwise}.
\end{cases} 
$$    
\end{lemma} 

\begin{lemma}
\label{generalcoeff}
For any $i,n,s$ with $1 \le i \le n$ and  
$0 \le s \le \lfloor \frac{n+1}{2} \rfloor$, we have 
\begin{multline*} 
\sum_{j=1}^n \sum_{k=0}^n b_{j,k} e_k(q(X_1),\ldots,q(X_n)) p(X_i)^s \left( p(X_i)^{j-(1+[j<2s])s}-q(X_i)^{j-(1+[j<2s])s} \right) \\
=\prod_{k=1}^n (p(X_i) - q(X_k)),
\end{multline*}                   
where 
$$
b_{j,k} = 
\begin{cases} 
(-1)^{k+1} \left( \binom{s+n-k-j}{2 (-2s+n-k)+j} 
+ \binom{s+n-k-j-1}{2 (-2s+n-k)+j-1} \right),
& n-2s+1 \le k \le n -s-1 \,  \&   \, j \le 2s-1, \\
(-1)^{k} [j+k=n+s], &  n - s-1 < k \, \&   \, j \le 2s-1, \\
(-1)^k [j+k=n], & \text{otherwise}. 
\end{cases} 
$$

\end{lemma} 

\begin{proof}
As in the proof of Lemma~\ref{coeffbase}, we expand the product on the right-hand side in the statement of the lemma, then apply 
the identity $e_k=e_k^{(i)}+q(X_i) e_{k-1}^{(i)}$ (using the same notation) and compare the coefficients of 
the $e^{(i)}_k$'s on both sides. Thus we need to show 
\begin{multline*} 
(-1)^{k} \left(p(X)^{n-k}-p(X_i)^{n-k-1} q(X) \right) \\
=
\sum_{j=1}^n (b_{j,k}+b_{j,k+1}q(X)) p(X)^{s} \left( p(X)^{j-(1+[j<2s])s}-q(X)^{j-(1+[j<2s])s} \right)
\end{multline*} 
for $k=0,1,\ldots,n-1$. It suffices to show 
$$ 
(-1)^{k} \left(p(X)^{n-k-s}-q(X)^{n-k-s} \right) \\
=
\sum_{j=1}^n b_{j,k} \left( p(X)^{j-(1+[j<2s])s}-q(X)^{j-(1+[j<2s])s} \right).
$$
The sum on the right-hand side is split into the range $1 \le j \le 2s-1$ and $2s \le j \le n$. 
As for the second part, note that 
$$
\sum_{j=2s}^n (-1)^k [j+k=n] \left( p(X)^{j-s}-q(X)^{j-s} \right)
= [n-k \ge 2s] (-1)^k \left( p(X)^{n-k-s}-q(X)^{n-k-s} \right),
$$
so that we need to show 
$$ 
[n-k \le 2s-1] (-1)^{k} \left(p(X)^{n-k-s}-q(X)^{n-k-s} \right) \\
=
\sum_{j=1}^{2s-1} b_{j,k} \left( p(X)^{j-2s}-q(X)^{j-2s} \right).
$$
Next we distinguish between the cases $k \le n-2s$, $n-2s+1 \le k \le n-s-1$ and $n-s \le k$.
If $k \le n-2s$, then the left-hand side vanishes. We have $b_{j,k}=(-1)^k [j+k=n]$ in this case, so that 
$b_{j,k} \not= 0$ only if $j=n-k$, but since $n-k \ge 2s$, this $j$ is not in the range. If 
$n-s \le k$, then the left hand side is 
$$
(-1)^{k} \left(p(X)^{n-k-s}-q(X)^{n-k-s} \right) 
$$ 
unless $s=0$, but this case is easy to check. Then we have $b_{j,k}=(-1)^k [j+k=n+s]$, so that the 
right-hand side is equal to 
$$
(-1)^k \left( p(X)^{n+s-k-2s}-q(X)^{n+s-k-2s} \right)
$$
which matches the left-hand side, 
since $1 \le n+s-k \le 2s-1$ unless $n+s-k=2s$. If $n+s-k=2s$ then the right-hand side is zero, but so is the left-hand side.

It remains to show 
\begin{multline*} 
\left(p(X)^{n-k-s}-q(X)^{n-k-s} \right) \\ = \sum_{j=1}^{2s-1}  \left( \binom{s+n-k-j}{2 (-2s+n-k)+j} 
+ \binom{s+n-k-j-1}{2 (-2s+n-k)+j-1} \right) \left( q(X)^{j-2s}-p(X)^{j-2s} \right)
\end{multline*} 
if $n-2s+1 \le k \le n-s-1$, which we will show for any integer $s$. For this, it suffices to show the following two identities.
\begin{multline*} 
\sum_{j=1}^{2s-1} \left( \binom{s+n-k-j}{2(-2s+n-k)+j} + \binom{s+n-k-j-1}{2(-2s+n-k)+j-1} \right) p(X)^{j-2s} \\ = 
(-1)^{n+s+k} p(X)^{n-k-s} X^{-3(n-k-s)} + p(X)^{-2n+2k+2s} X^{-3(k-n+s)} 
\end{multline*} 
\begin{multline*} 
\sum_{j=1}^{2s-1} \left( \binom{s+n-k-j}{2(-2s+n-k)+j} + \binom{s+n-k-j-1}{2(-2s+n-k)+j-1} \right) q(X)^{j-2s} \\= 
(-1)^{n+s+k} q(X)^{n-k-s} X^{3(n-k-s)} + q(X)^{-2n+2k+2s} X^{3(k-n+s)}
\end{multline*} 
To see this, we have used $q(X)=-p(X) X^{-3}$. As $q(X)=p(X^{-1})$, the two identities are equivalent and it suffices to 
show the first.

We perform the index transformation $i=2s-j$ in the first identity and obtain the following equivalent identity.
\begin{multline*} 
\sum_{i=1}^{2s-1} \left( \binom{-s+n-k+i}{2(-s+n-k)-i} + \binom{-s+n-k+i-1}{2(-s+n-k)-i-1} \right) p(X)^{-i} \\ = 
(-1)^{n+s+k} p(X)^{n-k-s} X^{-3(n-k-s)} + p(X)^{-2n+2k+2s} X^{-3(k-n+s)}
\end{multline*} 
The sum on the left-hand side does not change if we extend the summation over all $i \ge 1$. Also we set $d=-s+n-k$ and obtain 
the following equivalent identity 
$$
\sum_{i \ge 1} \left( \binom{d+i}{2d-i} + \binom{d+i-1}{2d-i-1} \right) p(X)^{-i} \\ = 
(-1)^{d} p(X)^{d} X^{-3 d} + p(X)^{-2d} X^{3d}
$$
This identity was proved in Lemma~\ref{coeffbase}. 
\end{proof} 

\end{document}
